\chapter{The Poisson Clock}

\begin{center}\itshape
Buses arrive at rate $\lambda$. The time you still have to wait is exponential of that rate,\\
so the mean wait is $1/\lambda$, not half of the mean gap.
\end{center}

An exponential clock has no memory. A Poisson process is a sequence of independent exponential gaps. Conditional on how many points fall in an interval, those points are uniform, and that one fact gives thinning, the inspection paradox, and the stationary queue.

\section{The exponential law}

A random variable $E$ is exponential of rate $\lambda>0$ when $P(E>t)=e^{-\lambda t}$ for $t\ge0$. It has density $\lambda e^{-\lambda t}$ on $(0,\infty)$ and mean $1/\lambda$.

\begin{proposition}[memorylessness]\label{prop:memoryless}
Let $T>0$ be a random variable with a continuous survival function $G(t)=P(T>t)$, not almost surely infinite, and with $G(s+t)=G(s)G(t)$ for all $s,t\ge0$. Then $T$ is exponential.
\end{proposition}
\begin{proof}
$G(0)=1$. If $G(1)=0$, then $G(1/n)^n=0$, so $G(1/n)=0$, and continuity forces $G(0)=0$, which is absurd. If $G(1)=1$, then $G(n)=1$ and $T$ is infinite. So $G(1)=e^{-\lambda}$ for some $\lambda>0$. Then $G(p/q)=e^{-\lambda p/q}$ on the positive rationals, and continuity extends the formula to every $t\ge0$.
\end{proof}

The lack of memory is the computation $P(E>s+t\mid E>s)=e^{-\lambda t}$. A bus whose gaps are exponential does not become more due because you have already waited.

\section{Arrival times}

\begin{lemma}[joint density of the partial sums]\label{lem:gamma}
Let $E_1,\ldots,E_n$ be independent exponential random variables of rate $\lambda$, and let $T_k=E_1+\cdots+E_k$. The vector $(T_1,\ldots,T_n)$ has density $\lambda^n e^{-\lambda t_n}$ on the simplex $0<t_1<\cdots<t_n$, and density $0$ elsewhere.
\end{lemma}
\begin{proof}
The independent gaps have joint density $\lambda^n\exp(-\lambda(e_1+\cdots+e_n))$ on the positive orthant. To compute the density of the partial sums without assuming a general multivariable change-of-variables theorem, start with a continuous test function $g$ supported in a compact simplex. Integrate successively in $e_n,e_{n-1},\ldots,e_1$. Holding the earlier gaps fixed, substitute $t_k=e_1+\cdots+e_k$ in the corresponding one-variable integral; its derivative with respect to $e_k$ is $1$. The new limits are $t_{k-1}<t_k$, and $e_1+\cdots+e_n=t_n$. The resulting iterated integral is
\[
 \int_{0<t_1<\cdots<t_n}g(t_1,\ldots,t_n)\lambda^n e^{-\lambda t_n}\,\dd t_n\cdots\dd t_1.
\]
This is successive one-variable substitution plus the compact-rectangle Fubini rule of Proposition~\ref{prop:fubini}; exponential tails extend it to the entire simplex. Hence the stated joint density. The triangular transformation has Jacobian determinant $1$, but that general terminology is not an extra prerequisite.
\end{proof}

\begin{definition}[Poisson process]\label{def:pp}
Let $E_i$ be i.i.d.\ exponential of rate $\lambda$, let $T_0=0$ and $T_n=E_1+\cdots+E_n$, and set $N(t)=\#\{n\ge1:T_n\le t\}$. This is the Poisson process of rate $\lambda$.
\end{definition}

The variables $E_i$ have moments of every order: $\mathbb E[E^k]=k!/\lambda^k$, by the integration by parts of Exercise~\ref{ex:gamma} after scaling. Theorem~\ref{thm:slln4} gives $T_n/n\to 1/\lambda$ almost surely, so $T_n\to\infty$ and $N(t)$ is finite.

\begin{theorem}[counts and conditional uniformity]\label{thm:uniform}
For each $t>0$ and integer $n\ge0$, $P(N(t)=n)=e^{-\lambda t}(\lambda t)^n/n!$. Conditional on $N(t)=n$, the arrival times $(T_1,\ldots,T_n)$ are distributed as the order statistics of $n$ independent uniforms on $[0,t]$.
\end{theorem}
\begin{proof}
The event $\{N(t)=n\}$ is $\{T_n\le t<T_n+E_{n+1}\}$. Independence of $E_{n+1}$ and Lemma~\ref{lem:gamma} give
\[
 P(N(t)=n,\,(T_1,\ldots,T_n)\in B)
 =\int_{B\cap\Delta_n(t)}\lambda^n e^{-\lambda t}\,\dd t_1\cdots\dd t_n,
\]
where $\Delta_n(t)$ is the simplex inside $[0,t]^n$. The volume of that simplex is $t^n/n!$, because the $n!$ possible orderings of $n$ distinct coordinates partition the cube up to boundaries of volume zero. Therefore $P(N(t)=n)=e^{-\lambda t}(\lambda t)^n/n!$, and the conditional distribution is uniform on the simplex, which is the law of the order statistics.
\end{proof}

\begin{lemma}[colouring]\label{lem:colouring}
Let $N$ be Poisson with mean $\mu$. Conditional on $N=n$, split the $n$ points into $k$ colours with probabilities $p_1,\ldots,p_k$. The colour counts are independent, and colour $j$ is Poisson with mean $\mu p_j$.
\end{lemma}
\begin{proof}
For $n_1+\cdots+n_k=n$,
\[
 P(N_1=n_1,\ldots,N_k=n_k)
 =e^{-\mu}\frac{\mu^n}{n!}\cdot\frac{n!}{n_1!\cdots n_k!}p_1^{n_1}\cdots p_k^{n_k}
 =\prod_{j=1}^k e^{-\mu p_j}\frac{(\mu p_j)^{n_j}}{n_j!},
\]
where $e^{-\mu}=\prod_j e^{-\mu p_j}$. A product of the marginal masses is independence.
\end{proof}

\begin{theorem}[independent increments]\label{thm:increments}
For $0=t_0<\cdots<t_m$, the increments $N(t_j)-N(t_{j-1})$ are independent, and $N(t_j)-N(t_{j-1})$ is Poisson of mean $\lambda(t_j-t_{j-1})$.
\end{theorem}
\begin{proof}
Conditional on $N(t_m)=n$, the points are uniform in $[0,t_m]$ by Theorem~\ref{thm:uniform}. Their counts in the subintervals are multinomial. Lemma~\ref{lem:colouring}, applied to the Poisson number $N(t_m)$, makes those counts independent Poissons. The mean in an interval of length $\ell$ is $\lambda t_m\cdot(\ell/t_m)=\lambda\ell$.
\end{proof}

\begin{figure}[ht]
\centering
\fig{ch10-poissonclock}
\caption{A schematic timeline of original arrival points (blue) and an independent selection kept after thinning (rust). The dividing mark $t$ separates two disjoint time intervals. Arrival counts on disjoint intervals are independent by Theorem~\ref{thm:increments}; independently colouring each point gives Theorem~\ref{thm:thinning}. Positions and colours here illustrate the mechanism, not simulated or observed data.}
\label{fig:poissonclock}
\end{figure}

\section{Thinning, inspection, and a queue}

\begin{theorem}[thinning]\label{thm:thinning}
Keep each point of a Poisson process of rate $\lambda$ independently with probability $p$. The kept points form a Poisson process of rate $\lambda p$, independent of the discarded points, which form a Poisson process of rate $\lambda(1-p)$.
\end{theorem}
\begin{proof}
On each finite interval the count is Poisson, and Lemma~\ref{lem:colouring} splits it. The increments on disjoint intervals remain independent because the original increments were independent and the coins are fresh.
\end{proof}

Superposition is the same lemma read backwards: independent Poisson processes of rates $\lambda$ and $\mu$ merge into one of rate $\lambda+\mu$.

Fix a time $t>0$ in the process that starts at $0$. The residual life is the wait $R_t=T_{N(t)+1}-t$ until the next point.

\begin{proposition}[residual life]\label{prop:residual}
For every $t\ge0$ and $s\ge0$, $P(R_t>s)=e^{-\lambda s}$. The mean wait is $1/\lambda$, whatever the clock already shows.
\end{proposition}
\begin{proof}
The event $\{R_t>s\}$ is the event of no point in $(t,t+s]$. Theorem~\ref{thm:increments} says that increment is Poisson of mean $\lambda s$, so its probability of equalling zero is $e^{-\lambda s}$. The mean of a nonnegative variable is the integral of its tail, $\int_0^\infty e^{-\lambda s}\,\dd s=1/\lambda$.
\end{proof}

The age is different if the process has a beginning. On the half-line the time since the previous point is at most $t$. To give the age the same law, run an independent Poisson process of rate $\lambda$ backward from $0$, and use its points as the past. The forward residual and the backward age are then independent exponential variables of rate $\lambda$, by the same increment computation on the two sides. Their sum is the gap that covers the instant $0$. Exercise~\ref{ex:twoexp} identifies the density of that sum as $\lambda^2 u e^{-\lambda u}$ on $(0,\infty)$, whose mean is $2/\lambda$. A typical gap, one of the original exponential variables, has mean $1/\lambda$. You are more likely to fall into a long gap, because a long gap covers more of the line. If buses arrive at rate $1/10$ per minute, the mean wait is ten minutes, and the mean length of the gap you stepped into is twenty minutes.

\begin{theorem}[Le Cam]\label{thm:lecam}
Let $X_i$ be independent Bernoulli random variables with means $p_i$, and let $S=\sum X_i$. Let $N$ be Poisson with mean $\sum p_i$. Then
\[
 \|\mathcal L(S)-\mathcal L(N)\|\le\sum_i p_i^2.
\]
\end{theorem}
\begin{proof}
We construct one exact coupling for each $i$, using independent auxiliary randomness across indices. Put $p=p_i\in[0,1]$ and draw $N_i\sim\operatorname{Poisson}(p)$. If $N_i\ge1$, let $X_i=1$. If $N_i=0$, let $X_i=1$ with conditional probability
\[
\alpha_p=\frac{p-(1-e^{-p})}{e^{-p}}.
\]
This is a valid probability: $e^{-p}\ge1-p$ implies $\alpha_p\ge0$, and $p\le1$ gives $\alpha_p\le1$. Then
\[
P(X_i=1)=P(N_i\ge1)+P(N_i=0)\alpha_p=(1-e^{-p})+p-1+e^{-p}=p,
\]
so $X_i$ really is Bernoulli$(p)$, unlike the indicator $\mathbf1_{\{N_i\ge1\}}$, whose success probability would be $1-e^{-p}$. The copies for distinct $i$ are independent by construction.

A discrepancy occurs only when $N_i=0$ but $X_i=1$, or when $N_i\ge2$ (since then $X_i=1\ne N_i$). Its probability is
\[
\begin{aligned}
P(X_i\ne N_i)&=e^{-p}\alpha_p+P(N_i\ge2)\\
&=p-1+e^{-p}+1-(1+p)e^{-p}
=p(1-e^{-p})\le p^2.
\end{aligned}
\]
The union bound gives $P(\sum_iX_i\ne\sum_iN_i)\le\sum_i p_i^2$. Total variation is at most the mismatch probability of any coupling, and the independent Poisson counts sum to a Poisson variable of parameter $\sum_i p_i$ by Chapter~3.
\end{proof}

\begin{theorem}[M/M/1]\label{thm:mm1}
Customers arrive as a Poisson process of rate $\lambda$ and require independent exponential service of rate $\mu>\lambda$. The number in the system is a Markov chain on $\{0,1,2,\ldots\}$. It has stationary distribution $\pi_k=(1-\rho)\rho^k$ with $\rho=\lambda/\mu$, and the stationary mean number in the system is $\rho/(1-\rho)$.
\end{theorem}
\begin{proof}
The chain moves from $k\ge1$ to $k+1$ at rate $\lambda$ and to $k-1$ at rate $\mu$, and from $0$ to $1$ at rate $\lambda$. Detailed balance $\pi_k\lambda=\pi_{k+1}\mu$ forces $\pi_{k+1}=\rho\pi_k$, so $\pi_k=\pi_0\rho^k$. The sum is $1$ when $\pi_0=1-\rho$. The mean is $\sum k\pi_k=\rho/(1-\rho)$. The balance equations determine the stationary law. The process itself is not a countable sample space: its exponential holding times are continuous. Construct the arrival and service clocks from independent uniform blocks as in Chapter~5. The total jump rate is at most $\lambda+\mu$, so in a finite interval the number of actual jumps is dominated by the finite number of arrivals of a Poisson clock of that rate; explosion cannot occur. Detailed balance for the rate matrix then makes the geometric distribution invariant. This establishes the stated stationary law, not a blanket theorem about all countable-state chains.
\end{proof}

For $\rho=0.9$ the mean number present is $9$. For $\rho=0.99$ it is $99$. The mean explodes as the server approaches saturation, although the chain remains positive recurrent for every $\rho<1$.

\section{Exercises}

\begin{exercise}[$\star$]\label{ex:twoexp}
Show that the sum of two independent exponential random variables of rate $\lambda$ has density $\lambda^2 t e^{-\lambda t}$ on $(0,\infty)$, and mean $2/\lambda$.
\end{exercise}

\begin{exercise}[$\star\star$]\label{ex:minexp}
Let $E_i$ be independent exponential random variables of rates $\lambda_i$. Show that $\min E_i$ is exponential of rate $\sum\lambda_i$, and that the probability the minimum is $E_j$ equals $\lambda_j/\sum\lambda_i$.
\end{exercise}

\begin{exercise}[$\star\star$]\label{ex:bus}
Buses follow a stationary Poisson process on the entire time line, of rate $1/10$ per minute (arrivals have occurred in the past as well as in the future). You arrive at a fixed time. What is your expected wait, and what is the expected length of the gap you occupy?
\end{exercise}

\begin{exercise}[$\star\star$]\label{ex:lecamnum}
Twenty independent events of probability $0.02$ occur or not. Bound the total variation distance between the number of occurrences and a Poisson random variable of mean $0.4$.
\end{exercise}

\begin{exercise}[$\star\star$]\label{ex:mm1mean}
In Theorem~\ref{thm:mm1}, derive $\sum_{k\ge0}k\rho^k=\rho/(1-\rho)^2$ for $|\rho|<1$, and conclude that the stationary mean is $\rho/(1-\rho)$.
\end{exercise}

\begin{exercise}[$\star$]\label{ex:superpose}
Use Lemma~\ref{lem:colouring} to show that the sum of independent Poisson random variables of means $\lambda$ and $\mu$ is Poisson of mean $\lambda+\mu$.
\end{exercise}

\section{Hints}

\begin{hint}{ex:twoexp}
Integrate the joint density on $e_1+e_2\le t$, or differentiate the distribution function.
\end{hint}
\begin{hint}{ex:minexp}
The minimum exceeds $t$ only if every variable does. The identity of the minimiser is a race of the hazards.
\end{hint}
\begin{hint}{ex:bus}
The residual life is a fresh exponential. The age of a bi-infinite process has the same law.
\end{hint}
\begin{hint}{ex:lecamnum}
Theorem~\ref{thm:lecam} with twenty copies of $p^2$.
\end{hint}
\begin{hint}{ex:mm1mean}
Differentiate $\sum\rho^k=1/(1-\rho)$.
\end{hint}
\begin{hint}{ex:superpose}
A Poisson$(\lambda+\mu)$ count, coloured with probabilities $\lambda/(\lambda+\mu)$ and $\mu/(\lambda+\mu)$, splits into the two summands.
\end{hint}

\section{Solutions}

\begin{solution}{ex:twoexp}
\writeup The joint density is $\lambda e^{-\lambda e_1}\lambda e^{-\lambda e_2}$ on the positive quadrant. The sum $T$ has
\[
 P(T>t)=\int_{e_1+e_2>t}\lambda^2 e^{-\lambda(e_1+e_2)}\,\dd e_1\,\dd e_2=e^{-\lambda t}(1+\lambda t),
\]
by integrating $e_1$ from $0$ to $\infty$ and $e_2$ from $(t-e_1)_+$ . Differentiating the distribution function $1-e^{-\lambda t}(1+\lambda t)$ produces the density $\lambda^2 t e^{-\lambda t}$. The mean is $2/\lambda$ by a further integration, or by adding the two means.
\end{solution}

\begin{solution}{ex:minexp}
\writeup $P(\min E_i>t)=\prod e^{-\lambda_i t}=e^{-(\sum\lambda_i)t}$. For $E_j$ to be the smallest at time $t$, it must fall in an infinitesimal interval there while all other clocks exceed $t$. Independence and the densities give the exact integral
\[
P(E_j=\min_i E_i)=\int_0^\infty\lambda_j e^{-\lambda_j t}\prod_{i\ne j}e^{-\lambda_i t}\,\dd t
=\frac{\lambda_j}{\sum_i\lambda_i}.
\]
Ties have probability zero for continuous independent clocks.
\end{solution}

\begin{solution}{ex:bus}
\writeup The forward wait is exponential of rate $1/10$, with mean $10$ minutes. Because the process is defined on the whole time line, the age since the previous arrival also has the exponential law with mean $10$ minutes (apply the no-point probability to the interval $[-t,0]$). Thus the expected length of the gap containing the arrival time is $10+10=20$ minutes. For a process started at time zero, the backward age is truncated and this conclusion would need a different model.
\end{solution}

\begin{solution}{ex:lecamnum}
\writeup Theorem~\ref{thm:lecam} bounds the distance by $20\cdot(0.02)^2=0.008$.
\end{solution}

\begin{solution}{ex:mm1mean}
\writeup Differentiate $\sum_{k\ge0}\rho^k=1/(1-\rho)$ to obtain $\sum_{k\ge1}k\rho^{k-1}=1/(1-\rho)^2$, hence $\sum k\rho^k=\rho/(1-\rho)^2$. Multiply by $1-\rho$.
\end{solution}

\begin{solution}{ex:superpose}
\writeup Let $N$ be Poisson$(\lambda+\mu)$ and colour each point independently red with probability $\lambda/(\lambda+\mu)$. Lemma~\ref{lem:colouring} says the red and blue counts are independent Poissons of means $\lambda$ and $\mu$, and their sum is $N$.
\end{solution}

\section*{Sources and further reading}

\begin{enumerate}
\item S.~D.~Poisson, \emph{Recherches sur la probabilit\'e des jugements}, Bachelier, Paris, 1837.
\item L.~Le Cam, An approximation theorem for the Poisson binomial distribution, \emph{Pacific J.\ Math.} \textbf{10} (1960), 1181--1197. Theorem~\ref{thm:lecam}, in the form proved above.
\end{enumerate}
